% Sources of the parts of the simulation model (design of Chapters 4 and 5). Written by hand, 19 Nov 2026.
\begin{table}[htbp]
\centering
\footnotesize
\setlength{\tabcolsep}{4pt}
\begin{tabularx}{\textwidth}{L{3.6cm}YL{2.6cm}}
\toprule
\textbf{Source} & \textbf{Parts of the simulation model} & \textbf{Used in} \\
\midrule
Learned from the training log & Net, routing trees, waiting-time trees and execution-time trees & All simulation runs \\
\addlinespace
Learned, but replaced in the conditional evaluation & Arrival tree, distribution of the case attributes and block-choice trees & E3 and the generative runs \\
\addlinespace
Taken from the observed log & Gate-In times and case attributes, block of each RMG stop with an observed counterpart, closed blocks from the yard utilisation data & E4 and E5 \\
\addlinespace
Fixed by a physical rule or by the terminal & Domain rules, stop guard for the planned number of stops, three lanes per block with the capacity checked at the start of a stop, allocation rule & E2 to E5 \\
\bottomrule
\end{tabularx}
\par\smallskip{\footnotesize\raggedright The learned trees are fitted by ProSiT. Where a part is corrected, the correction changes its training data or the mechanism that uses it, never the fitted values by hand (Chapter~\ref{chp:process_discovery_methodology}).\par}
\caption{What the simulation model learns, what is fixed and what is taken from the log.}
\label{tab:exp_sources}
\end{table}
